%=======================================================================
% Template para dissertação 
% Programa de Pós Graduação em Informática Aplicada da Universidade Federal 
% Rural de Pernambuco
%=======================================================================

%para incluir um arquivo por vez use \includeonly. 
%exemplos:
%\includeonly{capitulos/introducao,capitulos/bibliografia}
%\includeonly{capitulos/introducao,capitulos/conclusoes,capitulos/bibliografia}
%\includeonly{prefacios/capa2,prefacios/folha_rosto,prefacios/resumo,prefacios/abstract,capitulos/introducao,capitulos/bibliografia }

\documentclass [12pt,fleqn,oneside]{book}    % nomes e hifenação em português
\usepackage[brazil]{babel}         % Pacotes relativos à linguagem de entrada
\usepackage [utf8]{inputenc}     % aceitar texto acentuado em UTF-8
\usepackage {fancyhdr}             % cabeçalhos mais sofisticados
\usepackage {indentfirst}          % indentar 1a linha de cada seção

% Pacotes adicionados por mim
\usepackage{float}
\usepackage{nomencl}
\usepackage{multirow}
\usepackage{caption}
\usepackage{array}
\usepackage{tabularx}

% Pacotes relativos à inclusão de figuras
\usepackage {graphicx}         % para incluir figuras em PostScript
\usepackage{amsfonts,amsthm,amsopn,amssymb,latexsym}
\usepackage{geometry}
\usepackage[intlimits]{amsmath}
\usepackage{hyperref}
\usepackage{color}

% Dimensões da página
\usepackage{a4}                       % tamanho da página
\setlength{\textwidth}{16.5cm}        % largura do texto
\setlength{\textheight}{9.5in}        % tamanho do texto (sem head, etc)
\renewcommand{\baselinestretch}{1.4}  % espaçamento entre linhas
\addtolength{\topmargin}{-1cm}        % espaço entre o head e a margem
\setlength{\oddsidemargin}{-0.1cm}    % espaço entre o texto e a margem
\setlength{\parindent}{.0in}          % identação na 1a linha de cada parágrafo
\sloppy                               % Ser indulgente no preenchimento das linhas
\parskip 10pt                         % Espaçamento entre parágrafos

%alguns macros
\newtheorem{teo}{Teorema}[chapter]
\newtheorem{lema}[teo]{Lema}
\newtheorem{cor}[teo]{Corol\'{a}rio}
\newtheorem{prop}[teo]{Proposi\c{c}\~{a}o}
\newtheorem{defin}[teo]{Defini\c{c}\~ao}
\newtheorem{Af}[teo]{Afirma\c{c}\~{a}o}
\newtheorem{ex}[teo]{Exemplo}
\newtheorem{obs}[teo]{Observa\c{c}\~{a}o}
\newcommand{\R}{\ensuremath{\mathbb{R}}}
\newcommand{\Rn}{{\ensuremath{\mathbb{R}}}^{n}}
\newcommand{\Rm}{{\ensuremath{\mathbb{R}}}^{m}}
\newcommand{\Rnn}{{\ensuremath{\mathbb{R}}}^{ n \times n }}
\newcommand{\N}{\ensuremath{\mathbb{N}}}
\newcommand{\Z}{\ensuremath{\mathbb{Z}}}
\newcommand{\blackbox}{\rule{2mm}{2mm}}
\newcommand{\raiz}{\mbox{raiz}}
\newcommand{\contcaption}[1]{\vspace*{-0.1\baselineskip}\begin{center}#1\end{center}\vspace*{0.3\baselineskip}}

% Definindo as barras para a capa
\newcommand{\HRule}{\rule{\linewidth}{1mm}}

% Customização da Lista de Siglas
\makenomenclature
\renewcommand{\nomname}{Lista de Siglas}

% ====================================================================================


% Início do texto
\begin {document}

% Páginas sem numeração (capa, agradecimentos, etc)
\pagestyle {empty}

% Páginas iniciais
\begin{figure}[h]
\leavevmode
\begin{minipage}{\textwidth}
\includegraphics[scale=0.7]
{prefacios/capa/logo-ufrpe.eps}
\end{minipage}
\end{figure}
\vspace{-3.5cm}
{\bf
\begin{center}
{\Large
\hspace*{0cm}Universidade Federal Rural de Pernambuco \\
\hspace*{0cm}Departamento de Estatística e Informática}\\
\end{center}
}
\noindent
\begin{figure}[h]
\centering
\includegraphics[scale=0.5]{prefacios/capa/logo-bsi-presencial-v3-amp.eps}
\end{figure}


\vspace{2.5cm}
\noindent
\begin{center}
{\Large \bf Um Estudo sobre Lições Aprendidas na Adoção de Métodos Ágeis por Organizações de Software} \\[5cm]
{\Large Daniel Filype Silva Barreto}\\[6mm]
\end{center}


\vspace{1.5cm}
\begin{center}
{\large {\bf Recife}\\[6mm]
Março de 2014}
\end{center}
\newpage
           % capa ilustrativa
%
\hyphenation{}

\vspace*{0.0cm}
{\center
{\Large Daniel Filype Silva Barreto}\\[2.4cm]
{\huge \bf Um Estudo sobre Lições Aprendidas na Adoção de Métodos Ágeis por Organizações de Software}\\[2.0cm]
{\Large Orientador: Teresa Maciel}}\\[2.0cm]

{\raggedleft
\begin{minipage}[t]{8.3cm}
\setlength{\baselineskip}{0.25in}
Monografia apresentada ao Curso Bacharelado em Sistemas de Informação  da Universidade Federal Rural de Pernambuco, como requisito parcial para obtenção do título de Bacharel em Sistemas de Informação.\end{minipage}\\[2cm]}
\vspace{3cm}
{\center Recife \\[3mm]
Março de 2014 \\}

\newpage
\vspace*{18cm}
{\raggedleft
\begin{minipage}[t]{6.0cm}
\setlength{\baselineskip}{0.25in}
À minha mãe, Izabel\\
À minha namorada, Bruna\\
À minha orientadora, Teresa\\
Aos meus amigos\\
\end{minipage}\\[2cm]}



\newpage
\begin{center}
{\Large \bf Agradecimentos}
\end{center}
\vspace*{-0.06in}

À minha mãe, Izabel, que desde pequeno me mostrou o quanto a busca por conhecimento é importante.

À minha namorada, Bruna, por sempre estar ao meu lado em todos os momentos importantes da minha vida.

À minha orientadora, Teresa, que sempre esteve disponível e disposta a me ajudar.

Aos meus amigos de infância que, mesmo indiretamente, incentivaram-me a seguir em frente.


     % capa ABNT

% Numeração em romanos para páginas iniciais (sumários, listas, etc)
\pagenumbering {roman}
\pagestyle {plain}

% Geração do sumário, agradecimentos, e etc..
\chapter*{Resumo}
escreva o resumo do trabalho...

         % resumo em português
\chapter*{Abstract}

Agile development methodologies aim to build software in a faster and more efficient way, creating a minimum product that already brings value to the customer and, thenceforth, adding new features to it. This is a complete paradigm shift if we compare them with the planning-driven methodologies.

Changing drastically is never easy. Given this scenario, a research line focused on the study of this transition process from planning-driven methodologies to Agile has emerged. According to its studies, the main benefits from the adoption of agile methods are: customer satisfaction, more frequent deliveries, maintainability of the team morale high, improvement on the quality of the final product and so on. However, conquering all these benefits requires a lot of discipline and dedication.

Analyzing the results obtained with this research, a more common set of lessons learned was perceived. Among them, we can cite: ``customization and adaptability", ``experience, training and learning", ``engagement, commitment, discipline and teamwork" and ``technical and technological aspects".

\textbf{Keywords:} Agile methods, agile adoption, lessons learned

\tableofcontents                    % sumário
\listoftables                       % lista de tabelas
\listoffigures                      % lista de figuras
\printnomenclature		     % lista de siglas
%\listofprograms                    % lista de programas
\clearpage


%=======================================================================
% Definir um estilo para páginas completas (cabeçalho + rodapé)
\pagestyle {fancyplain}

% Marca de capítulo do tipo "2. Blábláblá..." no cabeçalho
\renewcommand{\chaptermark}[1]{\markboth{\thechapter.\ {#1}}{}}

% Para aumentar o tamanho da caixa para o cabeçalho
\addtolength{\headheight}{\baselineskip}

% Definindo o conteúdo do cabeçalho...
\fancyhf{}
\fancyhead[L,L]{\nouppercase{\textsf{\leftmark}}} %comentar
\fancyhead[R,R]{\thepage}

% Para garantir que a primeira página de cada capítulo...
\fancypagestyle{plain}{
\fancyhf{}                          % não terá head ou foot
\renewcommand{\headrulewidth}{0pt}} % nem linhas de cabeçalho
%=======================================================================


% Capítulos são numerados em algarismos arábicos
\setcounter{page}{0} \pagenumbering{arabic}

%incluir os arquivos dos capítulos
% Os arquivos de cada capítulo devem iniciar com \chapter{Nome do Capítulo}
% (o nome deve terminar com .tex)
\chapter{Introdução}

\section{Apresentação}

O Manifesto Ágil \cite{agileManifesto}, criado em 2001, revolucionou o mercado de Tecnologia da Informação, pois, até então, a maneira de se desenvolver software mostrou-se falha em diversos aspectos. Contudo, a transição entre os modelos antigos (conhecidos como “plan-driven”) e o moderno (conhecido como Ágil) é uma tarefa não-trivial.

As principais motivações por trás da adoção ágil de desenvolvimento de software são: a melhoria da qualidade do produto final, aumento da moral dos desenvolvedores e satisfação do cliente. Entretanto, adoção ágil sempre vem com desafios especiais e, para que ela ocorra com sucesso, mudanças fundamentais na cultura da organização são necessárias \cite{Hajjdiab2011}. Este projeto visa expor lições aprendidas por empresas de desenvolvimento de software que passaram por este processo.

\section{Justificativas}

Muitas abordagens diferentes podem ser aplicadas a desenvolvimento de software \cite{Kettunen2010}, cada uma com suas peculiaridades. De acordo com \cite{Shore2007}, metodologias denominadas ágeis tornaram-se bastante populares, sendo utilizadas por grandes corporações como: Google, Yahoo!, Symantec, Microsoft, etc. Entretanto, não existe bala de prata no que tange a desenvolvimento de software. Não se deve simplesmente aplicá-las porque são populares. Existem vários estudos na literatura que mostram casos falhos de adoção de métodos ágeis \cite{Krasteva2008}. É preciso entender a fundo estas novas metodologias, analisar seus pontos positivos e negativos e, principalmente, refletir se elas agregariam valor ao produto final a ser entregue.

Sendo assim, é de suma importância a coleta de experiências para que se possa aprender sobre esta transição com aqueles que já passaram por este processo e que, além disso, querem contribuir com a comunidade de desenvolvimento de software.

\section{Objetivos} 

\subsubsection{Objetivo Geral}

\begin{itemize}
	\item Disponibilizar para acesso da indústria e academia de software um conjunto de lições aprendidas mais relatadas com a adoção de desenvolvimento ágil por organizações de software.
\end{itemize}

\subsubsection{Objetivos Específicos}

\begin{itemize}
	\item Investigar, através de pesquisa exploratória, a necessidade e contribuição de trabalhos desta natureza.
	\item Adquirir conhecimento sobre fundamentos teóricos relacionados como desenvolvimento ágil de software e lean.
	\item Realizar uma revisão utilizando como base a metodologia de revisão sistemática \cite{Barbara04} para coletar lições aprendidas provenientes de trabalhos publicados pela comunidade científica.
	\item Realizar uma pesquisa exploratória para coletar lições aprendidas provenientes de trabalhos publicados pela indústria de software nas principais conferências neste contexto.
	\item Identificar um conjunto de benefícios e desafios mais citados a partir dos resultados das pesquisas realizadas.
	\item Validar o trabalho realizado através de uma metodologia de survey \cite{Babbie1990}, aplicando-se questionários com profissionais brasileiros reconhecidos na área, assim como em empresas que adotam o desenvolvimento ágil.
	\item Gerar um protótipo de um banco de dados de lições aprendidas para acesso livre pela comunidade de software.
\end{itemize}

\section{Contribuições obtidas}

Com este trabalho, foi possível criar uma fonte única de informações relacionadas a processos de adoção ágil em empresas de desenvolvimento de software. Além disso, foi possível iniciar da criação de um protótipo de banco de dados de lições aprendidas que, dependendo do interesse da comunidade, pode vir a se tornar uma grande ferramenta colaborativa.

\section{Organização do trabalho}

Este trabalho é composto por mais três capítulos. No capítulo 2 está exposto o método utilizado para revisar a Literatura, juntamente com o conjunto de artigos científicos e relatos de experiência selecionados. No capítulo 3 encontra-se o resultado do mapeamento das lições aprendidas relatadas pelos trabalhos selecionados. O capítulo 4 apresenta a validação do que foi construído com o projeto (através de um questionário), suas conclusões e possíveis futuros trabalhos.


\chapter{Revisão da Literatura}
	Serão apresentados neste capítulo os seguintes pontos: a estratégia de pesquisa adotada neste projeto de conclusão de curso, os desafios e as lições aprendidas documentadas por organizações de desenvolvimento de software que tentaram (com sucesso ou não) adotar ágil.
	\section{Método de pesquisa}
		O processo de pesquisa utilizado neste projeto foi baseado no método chamado revisão sistemática. Não foi necessário ser completamente fidedigno ao método dado a restrições de tempo, contingente (este foi um trabalho desenvolvido individualmente) e nível de detalhe.
		
		Segundo Barbara Kitchenham \cite{Barbara04}, uma revisão sistemática da literatura consiste em identificar, avaliar e interpretar todas pesquisas disponíveis relevantes a uma determinada questão, tópico ou área de interesse. Este tipo de pesquisa requer a definição de alguns pontos:
		\begin{itemize}
			\item A(s) questão(ões) a ser(em) respondida(s) ao final da pesquisa
			\item A estratégia utilizada (fontes consultadas, termos buscados, etc.)
			\item Critérios de seleção e exclusão de artigos
		\end{itemize}
		A necessidade da adoção de uma revisão sistemática decorre da exigência de se sumarizar todo um conjunto de informações existentes no que tange à adoção de metodologias ágeis por parte de empresas de desenvolvimento de software. A partir deste estudo, pode-se criar um arcabouço para a elaboração dos questionários que foram utilizados neste projeto.
		\subsection{Definição da pergunta}
			O primeiro passo foi definir a questão a ser respondida ao final da pesquisa: ``\textit{Quais os desafios e benefícios da adoção de metodologias ágeis por parte de empresas de desenvolvimento de software?}"
		\subsection{Bases de dados relevantes}
			Após a definição da pergunta a ser respondida, o próximo passo foi definir as fontes de informação mais confiáveis e relevantes para o tema (Tabela \ref{tab:basesDeDados}).
			\begin{table}[H]
				\centering
				\begin{tabular}{| l | r |} \hline \textbf{Nome} & \textbf{Referência} \\ \hline
					CAPES & http://www.capes.gov.br/ \\ \hline
					ACM & http://www.acm.org/ \\ \hline
					IEEE & http://ieeexplore.ieee.org/ \\ \hline
					Google Scholar & http://scholar.google.com/ \\ \hline
				\end{tabular}
				\caption{Resumo das bases de dados utilizadas na pesquisa}
				\label{tab:basesDeDados}
			\end{table}
			\nomenclature{CAPES}{Coordenação de Aperfeiçoamento de Pessoal de Nível Superior}%
			\nomenclature{ACM}{Association for Computing Machinery}%
			\nomenclature{IEEE}{Instituto de Engenheiros Eletricistas e Eletrônicos}%
		\subsection{Palavras-chave}
			Neste passo, foram definidas as palavras-chave (Tabela \ref{tab:palavrasChave}) buscadas nas bases de dados selecionadas.
			\begin{table}[H]
				\centering
				\begin{tabular}{| l |} \hline \textbf{Palavras-chave} \\ \hline
					Agile Adoption \\ \hline
				\end{tabular}
				\caption{Conjunto de palavras-chave utilizadas na pesquisa}
				\label{tab:palavrasChave}
			\end{table}
		\subsection{Critérios de exclusão}
			//Área para definir critérios de descarte após a pesquisa ter sido feita (colocar números)
			\begin{table}[H]
				\centering
				\begin{tabular}{| c | l | r |} \hline \textbf{Critério de exclusão} & \textbf{Base de dados}  & \textbf{Quantidade} \\ \hline
					\multirow{4}{*}{-}
						& ACM & 1 \\ \cline{2-3}
						& CAPES & 1 \\ \cline{2-3}
						& Google Scholar & 1 \\ \cline{2-3}
						& IEEE & 1 \\ \cline{2-3}
					\hline \hline
					\multirow{4}{*}{Artigos entre 2010 e 2013} 
						& ACM & 1 \\ \cline{2-3}
						& CAPES & 1 \\ \cline{2-3}
						& Google Scholar & 1 \\ \cline{2-3}
						& IEEE & 1 \\ \cline{2-3}
					\hline \hline
					\multirow{4}{*}{Palavras-chave no título e/ou resumo} 
						& ACM & 1 \\ \cline{2-3}
						& CAPES & 1 \\ \cline{2-3}
						& Google Scholar & 1 \\ \cline{2-3}
						& IEEE & 1 \\ \cline{2-3}
					\hline \hline
					\multirow{4}{*}{Análise crítica}
						& ACM & 1 \\ \cline{2-3}
						& CAPES & 1 \\ \cline{2-3}
						& Google Scholar & 1 \\ \cline{2-3}
						& IEEE & 1 \\ \cline{2-3}
					\hline
				\end{tabular}
				\caption{Quantidade de materiais encontrados em cada passo da pesquisa}
				\label{tab:quantidadeDeMateriais}
			\end{table}

\chapter{Análise qualitativa de lições aprendidas}

Após a seleção do conjunto de trabalhos primários desta pesquisa (16 artigos científicos e 14 relatos de experiência), as lições aprendidas reportadas por eles foram devidamente mapeadas e categorizadas. Neste capítulo encontra-se o resultado deste mapeamento.

\section{Mapeamento e categorização}
Para obter-se uma maior visibilidade do que foi encontrado, as diversas lições aprendidas reportadas pela base da literatura científica e relatos de experiência foram multidimensionalmente mapeadas.

Foi considerada a existência real de uma categoria quando pelo menos 6 trabalhos (20\% do material analisado) referenciavam alguma lição aprendida que se encaixasse naquela categoria (Tabela \ref{tab:mapeamentoCategorias}).

\begin{table}[H]
	\centering
	\begin{tabularx}{\linewidth}{ | p{6cm} | X | X | }
		\hline 
		\textbf{Categorias} & \textbf{Artigos científicos} & \textbf{Relatos de experiência} \\ 
		\hline 
		Experiência, treinamento e aprendizado & \cite{Hajjdiab2011}, \cite{Block2011}, \cite{Adobe2012}, \cite{Cisco2011}, \cite{Lapham2012}, \cite{Eunha2012}, \cite{Claudia2013}, \cite{Asnawi2012}, \cite{Fitzgerald2013} & \cite{Stefano2013}, \cite{Rodrigues2013}, \cite{Bastos2013}, \cite{Maciel2013}, \cite{Karaj2013}, \cite{Piegas2012}, \cite{Vieira2013} \\ 
		\hline 
		Planejamento e gerenciamento de backlog & \cite{Hajjdiab2011}, \cite{Fitzgerald2013}, \cite{Block2011}, \cite{Adobe2012}, \cite{Bustard2013}, \cite{Korhonen2010}, \cite{Claudia2013} & \cite{Piegas2012}, \cite{Hui2013}, \cite{Parzinello2012} \\ 
		\hline 
		Apoio gerencial e dos clientes & \cite{Hajjdiab2011}, \cite{Cisco2011}, \cite{Claudia2013}, \cite{Arikpo2011} & \cite{Parzinello2012}, \cite{Stefano2013}, \cite{Bastos2013}, \cite{Maciel2013}, \cite{Srinath2012}, \cite{Piegas2012} \\ 
		\hline 
		Customização e adaptabilidade & \cite{Hajjdiab2011}, \cite{Block2011}, \cite{Asnawi2012}, \cite{Fitzgerald2013}, \cite{Bustard2013}, \cite{Microsoft2013}, \cite{Lapham2012}, \cite{Claudia2013}, \cite{Nokia2013} & \cite{Piegas2012}, \cite{Hui2013}, \cite{Rodrigues2013}, \cite{Bastos2013}, \cite{Maciel2013}, \cite{Ahmed2008}, \cite{Sahota2012}, \cite{Vieira2013} \\ 
		\hline 
		Confiança do time & \cite{Block2011}, \cite{Asnawi2012}, \cite{Claudia2013}, \cite{Nokia2013} & \cite{Parzinello2012}, \cite{Ahmed2008}, \cite{Piegas2012}, \cite{Bastos2013} \\ 
		\hline 
		Engajamento, comprometimento, disciplina e trabalho em equipe & \cite{Block2011}, \cite{Asnawi2012}, \cite{Lapham2012}, \cite{Microsoft2013}, \cite{Claudia2013}, \cite{Nokia2013}, \cite{Adobe2012}, \cite{Fitzgerald2013} & \cite{Piegas2012}, \cite{Parzinello2012}, \cite{Stefano2013}, \cite{Rodrigues2013}, \cite{Maciel2013}, \cite{Queiroz2013}, \cite{Bastos2013}, \cite{Ahmed2008} \\ 
		\hline 
		Aspectos técnicos e tecnológicos & \cite{Block2011}, \cite{Microsoft2013}, \cite{Korhonen2010}, \cite{Cisco2011}, \cite{Lapham2012}, \cite{Eunha2012}, \cite{Fitzgerald2013}, \cite{Arikpo2011}, \cite{Bustard2013}, \cite{Radha2012}, \cite{Nokia2013} & \cite{Piegas2012}, \cite{Queiroz2013}, \cite{Stefano2013}, \cite{Karaj2013} \\ 
		\hline 
		Compartilhamento de conhecimento & \cite{Asnawi2012}, \cite{Cisco2011}, \cite{Lapham2012}, \cite{Radha2012}, \cite{Eunha2012}, \cite{Ericsson2013} & \cite{Valerio2013}, \cite{Vieira2013}, \cite{Queiroz2013}, \cite{Bastos2013}, \cite{Maciel2013} \\ 
		\hline 
		Velocidade de entrega e produtividade & \cite{Adobe2012}, \cite{Fitzgerald2013}, \cite{Microsoft2013}, \cite{Cisco2011}, \cite{Korhonen2010}, \cite{Eunha2012}, \cite{Claudia2013} & \cite{Stefano2013}, \cite{Queiroz2013}, \cite{Maciel2013}, \cite{Hui2013}, \cite{Ahmed2008}, \cite{Piegas2012} \\ 
		\hline 
		Qualidade do produto final & \cite{Adobe2012}, \cite{Fitzgerald2013}, \cite{Bustard2013}, \cite{Lapham2012}, \cite{Eunha2012}, \cite{Claudia2013}, \cite{Korhonen2010} & \cite{Parzinello2012}, \cite{Maciel2013}, \cite{Ahmed2008} \\ 
		\hline 
		Tamanho da organização & \cite{Bustard2013}, \cite{Microsoft2013}, \cite{Claudia2013}, \cite{Korhonen2010}, \cite{Ericsson2013} & \cite{Maciel2013} \\ 
		\hline 
		Quebra de paradigma & \cite{Hajjdiab2011}, \cite{Block2011}, \cite{Korhonen2010}, \cite{Lapham2012}, \cite{Arikpo2011} & \cite{Stefano2013}, \cite{Bastos2013}, \cite{Maciel2013}, \cite{Parzinello2012}, \cite{Hui2013}, \cite{Ahmed2008}, \cite{Sahota2012} \\ 
		\hline 
		Comunicação remota & \cite{Adobe2012}, \cite{Microsoft2013}, \cite{Korhonen2010}, \cite{Radha2012} & \cite{Rodrigues2013}, \cite{Vieira2013}, \cite{Bastos2013}, \cite{Maciel2013} \\ 
		\hline 
		Cultura organizacional & \cite{Bustard2013}, \cite{Microsoft2013}, \cite{Claudia2013}, \cite{Nokia2013}, \cite{Eunha2012}, \cite{Fitzgerald2013} & \cite{Rodrigues2013}, \cite{Bastos2013}, \cite{Srinath2012}, \cite{Maciel2013}, \cite{Sahota2012} \\ 
		\hline 
	\end{tabularx}
	\caption{Mapeamento de categorias de lições aprendidas e suas respectivas referências}
	\label{tab:mapeamentoCategorias}
\end{table}

\section{Lições aprendidas}

\subsection{Experiência, treinamento e aprendizado}

\begin{table}[H]
	\centering
	\begin{tabularx}{\linewidth}{ | X | p{5cm} | } \hline \textbf{Lições aprendidas} & \textbf{Referências} \\ \hline
		É muito difícil aventurar-se em ágil sem um Agile Coach & \cite{Hajjdiab2011}, \cite{Block2011}, \cite{Adobe2012}, \cite{Cisco2011}, \cite{Lapham2012}, \cite{Eunha2012}, \cite{Claudia2013}, \cite{Stefano2013}, \cite{Rodrigues2013}, \cite{Bastos2013}, \cite{Maciel2013}, \cite{Karaj2013} \\ \hline
		Vivenciar um projeto-piloto é uma prática muito eficiente para se adquirir experiência & \cite{Hajjdiab2011}, \cite{Cisco2011}, \cite{Rodrigues2013}, \cite{Maciel2013} \\ \hline
		É preciso prática, não apenas estudo & \cite{Hajjdiab2011}, \cite{Asnawi2012}, \cite{Claudia2013}, \cite{Piegas2012}, \cite{Vieira2013} \\ \hline
		Trabalhar com equipes não-agéis pode atrapalhar o andamento do projeto & \cite{Adobe2012}, \cite{Maciel2013} \\ \hline
	\end{tabularx}
\end{table}

\subsection{Planejamento e gerenciamento de backlog}

\begin{table}[H]
	\centering
	\begin{tabularx}{\linewidth}{ | X | p{5cm} | } \hline \textbf{Lições aprendidas} & \textbf{Referências} \\ \hline
		Template & \cite{} \\ \hline
	\end{tabularx}
\end{table}

\subsection{Apoio gerencial e dos clientes}

\begin{table}[H]
	\centering
	\begin{tabularx}{\linewidth}{ | X | p{5cm} | } \hline \textbf{Lições aprendidas} & \textbf{Referências} \\ \hline
		Template & \cite{} \\ \hline
	\end{tabularx}
\end{table}

\subsection{Customização e adaptabilidade}

\begin{table}[H]
	\centering
	\begin{tabularx}{\linewidth}{ | X | p{5cm} | } \hline \textbf{Lições aprendidas} & \textbf{Referências} \\ \hline
		Template & \cite{} \\ \hline
	\end{tabularx}
\end{table}

\subsection{Confiança do time}

\begin{table}[H]
	\centering
	\begin{tabularx}{\linewidth}{ | X | p{5cm} | } \hline \textbf{Lições aprendidas} & \textbf{Referências} \\ \hline
		Template & \cite{} \\ \hline
	\end{tabularx}
\end{table}

\subsection{Engajamento, comprometimento, disciplina e trabalho em equipe}

\begin{table}[H]
	\centering
	\begin{tabularx}{\linewidth}{ | X | p{5cm} | } \hline \textbf{Lições aprendidas} & \textbf{Referências} \\ \hline
		Template & \cite{} \\ \hline
	\end{tabularx}
\end{table}

\subsection{Aspectos técnicos e tecnológicos}

\begin{table}[H]
	\centering
	\begin{tabularx}{\linewidth}{ | X | p{5cm} | } \hline \textbf{Lições aprendidas} & \textbf{Referências} \\ \hline
		Template & \cite{} \\ \hline
	\end{tabularx}
\end{table}

\subsection{Compartilhamento de conhecimento}

\begin{table}[H]
	\centering
	\begin{tabularx}{\linewidth}{ | X | p{5cm} | } \hline \textbf{Lições aprendidas} & \textbf{Referências} \\ \hline
		Template & \cite{} \\ \hline
	\end{tabularx}
\end{table}

\subsection{Velocidade de entrega e produtividade}

\begin{table}[H]
	\centering
	\begin{tabularx}{\linewidth}{ | X | p{5cm} | } \hline \textbf{Lições aprendidas} & \textbf{Referências} \\ \hline
		Template & \cite{} \\ \hline
	\end{tabularx}
\end{table}

\subsection{Qualidade do produto final}

\begin{table}[H]
	\centering
	\begin{tabularx}{\linewidth}{ | X | p{5cm} | } \hline \textbf{Lições aprendidas} & \textbf{Referências} \\ \hline
		Template & \cite{} \\ \hline
	\end{tabularx}
\end{table}

\subsection{Tamanho da organização}

\begin{table}[H]
	\centering
	\begin{tabularx}{\linewidth}{ | X | p{5cm} | } \hline \textbf{Lições aprendidas} & \textbf{Referências} \\ \hline
		Template & \cite{} \\ \hline
	\end{tabularx}
\end{table}

\subsection{Quebra de paradigma}

\begin{table}[H]
	\centering
	\begin{tabularx}{\linewidth}{ | X | p{5cm} | } \hline \textbf{Lições aprendidas} & \textbf{Referências} \\ \hline
		Template & \cite{} \\ \hline
	\end{tabularx}
\end{table}

\subsection{Comunicação remota}

\begin{table}[H]
	\centering
	\begin{tabularx}{\linewidth}{ | X | p{5cm} | } \hline \textbf{Lições aprendidas} & \textbf{Referências} \\ \hline
		Template & \cite{} \\ \hline
	\end{tabularx}
\end{table}

\subsection{Cultura organizacional}

\begin{table}[H]
	\centering
	\begin{tabularx}{\linewidth}{ | X | p{5cm} | } \hline \textbf{Lições aprendidas} & \textbf{Referências} \\ \hline
		Template & \cite{} \\ \hline
	\end{tabularx}
\end{table}

\chapter{Análise quantitativa de lições aprendidas}

%Incluindo referências bibliográficas
\bibliographystyle{plain} %define o estilo
\bibliography{./bibliografia/referencias} %busca o arquivo referencias.bib

% Incluindo anexos numerados com letras maiusculas.

%\appendix 
% No arquivo, começar com \chapter{Nome do Anexo}
%\include{posfacios/anexo1}

% Fim do texto
\end{document}
%=======================================================================
